% Consolidated reply to all reviewers, NeurIPS 2026 submission 28595.
% Single working document: R1 and R2 evolve together here.
%
% Sources for every number: neurips_review/results_table.csv (new runs) and
% figures/flyvis/*_fullsample.tex (full-sample twins of Tabs. 1, 2, Supp. 4, 7).
%
% OUTSTANDING (see neurips_review/rebuttal_QA.md and PLAN_R2.md):
%   AC   - Appendix C response (kernel -> ker_eps, C.1/C.2 split, noise sweep)
%   R1   - done (200 ms cadence run stopped; sweep now ends at 100 ms = 5 tau)
%   R2   - all four questions; every [pending] marks a number that does not exist
%   PVbi - not drafted
\documentclass[11pt]{article}
\usepackage[margin=1in]{geometry}
\usepackage{amsmath,amssymb}
\usepackage{booktabs}
\usepackage[hybrid]{markdown}
\usepackage[colorlinks=true,allcolors=blue]{hyperref}
\setlength{\parskip}{5pt}
\setlength{\parindent}{0pt}
\newcommand{\Rw}{R^2_{\widehat W}}
\newcommand{\Rt}{R^2_{\hat\tau}}
\newcommand{\Rv}{R^2_{\hat V^{\mathrm{rest}}}}
\newcommand{\Mat}[1]{\mathbf{#1}}
% Table colouring, as in neurips.tex: good > 0.9, bad < 0.3.
\usepackage{xcolor}
\newcommand{\good}[1]{\textcolor{green!50!black}{#1}}
\newcommand{\bad}[1]{\textcolor{orange}{#1}}
\newcommand{\tbd}{\multicolumn{1}{c}{\textcolor{gray}{\textemdash}}}

\title{}
\date{}
\author{}

\begin{document}
\maketitle
\vspace{-3.5em}

\section{Responses}

\subsection{AC comment}

% Intro paragraph

We thank the reviewers for their close reading of our work and their thoughtful feedback on both the strengths and weaknesses of our work.

\subsubsection{Discussion of degeneracy in Appendix C}

% COPY PASTE THIS INTO OPEN REVIEW.
% NOTE THAT WE ARE WRITING MARKDOWN INSIDE OVERLEAF HERE.
% This is a high-level rebuttal of the AC's concerns about Appendix C.
% TODO (Stephan, Cedric): ready for review
\begin{markdown}
AC:
> Reconcile the approximately 6.4\% null plus 4.0\% sloppy edge dimensions with the claimed 71\% exact null space. Establish mathematically why time-shifted trajectories yield exact rank-one blocks, or acknowledge that Eq. (11) describes an approximate/sloppy subspace and revise the associated conclusions.
> ... Highly correlated or time-shifted signals do not generally constitute identical columns and therefore do not establish an exact (k-1)-dimensional kernel.
> ...The paper consequently appears to conflate exact non-identifiability, numerical sloppiness, and approximate dynamical equivalence.
> This issue directly affects one of the paper’s central analytical contributions and its interpretation of why process noise improves recovery.

Eq. (11) describes an _approximate_ or _sloppy_ subspace and not an exact null space. We acknowledge that the language in Appendix C is inaccurate and we propose to change/augment Appendix C to:
* add a paragraph to clarify the different concepts of exact/approximate degeneracies of the inverse problem and parameter perturbations that yield exactly/approximately equal activity traces (dynamical equivalence).
* add a paragraph to explain that exact dynamical equivalence = exact degeneracy of the inverse problem for FlyVis ODE. This was omitted from the original draft and its inclusion will help the reader understand this section better. Specifically, it allows us to use the Gram matrix in the least squares formulation to find directions in parameter space where perturbations would leave the activity traces approximately unchanged.
* include a figure showing the spectrum of the _normalized singular values_ (defined as $\lambda / \lambda_{max}$) of the Gram matrix in the least squares analysis. We have defined exact and sloppy directions using thresholds of `1e-22` and `1e-12`. It is the case that approximately 70\% of the normalized singular values are less than `1e-6`.
* clarify that the subsection that explicitly constructs parameter space directions using the symmetries of the connectome is not an exact null space, but an approximate one. Our goal is to provide the reader with some intuition about _why_ a large null space arises with a biological connectome and naturalistic input.

We are not aware of publications that have leveraged the linearity of the commonly used voltage-based firing rate model to quantify the degeneracy of the inverse problem using the singular values of the Gram matrix. We assert this to be a novel contribution and a connection that has not been made prior.

\end{markdown}

\subsection{1bPK Q1}

- good point
- delta t mismatch series done, robust results
- simulation with another equation (see below)
- model mismatch is also setting non-zero synapses to zero
- conductance based synapses would require a simulation with such a model (Janne working on it), then experiments iwht the following GNNs (current -> should not work well, because mismatch, GNN with $g(v_i,v_j,a_i,a_j)$ should work, then we have to test this also against the current model which is simpler, but could be harder for that GNN to fit, because underconstrained?)
  - conductance based synapse simulation dows not make sense, because the simulation model has not been tuned to generate useful task outputs like the original FlyVis model

\section{Notes}

The present document is a point-to-point list of answers to AC concerns, to reviewers questions and weak points. 

For a few answers, I ran new preliminary experiments following the suggestions of the AC/reviewers:

% Re: Appendix C
We agree with the reviewers, and the AC's summary that the discussion around degeneracy could be further clarified. We assert that the analysis in our minds is accurate but the quality of presentation can be improved so it would be clear in the mind of a reader.

> The paper consequently appears to conflate exact non-identifiability, numerical sloppiness, and approximate dynamical equivalence.

Using the notation of Appendix C, supplementary eq. (4) and (7),
Non-identifiability: when the matrix $A_i$ has a non-empty kernel, or equivalently, when the singular values of the Gram matrix $A^T A$ are zero, then the inverse problem (least squares) has an exact degenerate subspace.
Sloppiness: when the normalized singular values $\lambda / \lambda_{max}$ are small compared to an arbitrary but reasonably chosen threshold (1e-12) the curvature of the Hessian (equal to the Gram matrix) along the corresponding singular vector is too low to accurately determine the parameters $\theta_i$.
Dynamical equivalence: if the parameters of the ODE are perturbed along a direction $\theta_i \rightarrow \theta_i + \eta_i$, while keeping the initial condition $v_i(0)$ constant, the trajectory $v_i(t)$ is invariant.

% Dynamical equivalence = Non-identifiability

For this special class of ODEs where we generate the solution trajectories by Euler integration from a fixed initial condition, non-identifiability and dynamical equivalence, are equivalent notions. To see this, let's consider the one step update between $v_i(0)$ and $v_i(1)$. This corresponds to the first row of $A_i$, and if we modified $\theta_i \rightarrow \theta_i + \eta_i$, where $\eta_i$ is in the kernel of $A_i$, the new solution generated would be identical to the original. Will augment the appendix with a note that explicitly defines these notions and points out this equivalence.

Re: supplementary eq (11) and associated subsection of appendix C. We assert approximate dynamical equivalence, not an exact degenerate null space that comprises 70 percent of the weights. We acknowledge that the text inaccurately states that this is an exact null space in some places. We also agree with the AC that "Highly correlated or time-shifted signals do not generally constitute identical columns and therefore do not establish an exact (k-1)-dimensional kernel.".

We defined the sloppy eigenspace using a threshold of 1e-12 on the normalized singular values $\lambda/\lambda_{max}$. However, the full spectrum shows that approximately 70\% of the normalized singular values are below 1e-6. The goal of the exact degeneracy and sloppy direction analysis is to explain why the Known ODE oracle (and the GNN) fails to capture the true parameters, since a reader may attribute this to a poor learning algorithm. Directions with normalized singular values between 1e-12 and 1e-6 are correctly recovered by least squares and the Known ODE. However, if we perturb the parameters along one of these directions, and perform Euler integration from the initial condition, we observe that the activity traces are highly similar to the original (rollout Pearson $r > 0.94$). The discussion around time-shifted signals is meant to provide the reader with some intuition about this approximately degenerate parameter space and we think is valuable. 

By demonstrating that highly discordant weight matrices produce accurate roll outs we make clear to the reader why the baseline models are able to predict the rollout so accurately while failing to capture the underlying mechanism as illustrated by the Jacobian analysis.

% See Vijay's slack message: https://comp-theory.slack.com/archives/C0AMSUXMBEX/p1785185308957079?thread_ts=1784917005.150139&cid=C0AMSUXMBEX
% Let's not rush to make an argument about noise lifting degeneracy.

% Cedric's original text
%- concern of AC about Appendix C, apparent discrepancies between the sloppy $6\%$ in identifiably and $70\%$, I show that the two values correspond to two different \hyperref[link1]{tolerances} with a metrics build on the ratio of eigen values. 
%Further I ran a new set of \hyperref[link2]{SVD analysis} as a function of model noise level. This was not asked, I wanted to demonstrate that this break the columnar degeneracies, it is not conclusive yet.

- R1 Q1, about model class match (GNN too close to simulation). I ran two new experiments, i.e. adding a slow adaptation \hyperref[link3]{current} to break the first-order ODE and \hyperref[link4]{temporal mismatch} between simulation and GNN training.

- R1 Q2, I reanalyzed all experiments to provide \hyperref[link4]{full-sample} $R^2$.

- R2 Q2, I trained GNN and Know-ODE on \hyperref[link5]{Calcium-like} observation model with different level of model noise and measurement noise. Before applying the model training, the Calcium data is deconvolved. The inverse is a regularized Tikhonov in Fourier space. Preliminary results are interesting, GNN recovers connectivity if measurement noise is zero, works needs to be done to handle measurement noise (again). This is a new result, so far I never managed to have a GNN working on Calcium data. I was taking other route than regularization, e.g. training directly on Calcium. I tried the deconvolution before but with ill regularization (see \href{https://github.com/.../calcium_note.pdf}{calcium note})  

- R2 Q3, I ran different test on \hyperref[link6]{out of distribution data}: conclusive.

- R2 Q4, I ran the \hyperref[link7]{Jacobians analysis} on GNN and Known ODE to complement published MLP results: conclusive  

- R3 Q1, I ran a \hyperref[link8]{strongly recurrent benchmark} with a Drosophila complex system connectome (156 cells, 10,263 edges, density 0.422): ok if model noise $\sigma >0$
\\
\newline
\newline

\section{Confidential note to the AC}

We cite prior arXiv preprints that share authors with this submission. Reviewers have treated them as establishing that our contribution is not novel. May we disclose this relationship to the Area Chair, and through which channel?

\section{AC Could the author response significantly change the recommendation}

\subsection{}

\subsection{}

\subsection{}

\subsection{Temporal and dynamical model mismatch}
See Reviewer 1bP answer to Q1. How much of "recovery" depends on (a) the GNN's hypothesis class containing the true generator and (b) $\Delta t_{inference} = \Delta t_{simulation}$ ? 




\subsection{}


\section{Comparison with other papers:}

\hyperref[link8]{\textbf{Allier et al.\ [2026]}}

\hyperref[link9]{\textbf{Shiu et al. [2024]}}

\hyperref[link10]{\textbf{BrainTrace et al.[2026]}}

\hyperref[link11]{\textbf{AMAG [2023]}}

\hyperref[link12]{\textbf{Yoon [2025]}}

\hyperref[link13]{\textbf{Das Fiete [2020]}}


\newpage
\section*{Area Chair 4ZQS [DEPRECATED] }
We thank the Area Chair for a meta-review detailed enough that each concern could
be tested rather than guessed at. We answer the six weaknesses in order.
Concerns~1 to~4 carry new measurements, run for this reply; $5$ and~$6$ we accept
in full. Throughout we state what we concede and what we will change.

\paragraph{What the concessions do not touch.}
They concern framing, the interpretation of Appx.~C, and the strength of the
claims. The measurements stand as reported. Unaffected: the
prediction/mechanism dissociation, three high-capacity predictors reaching
rollout $r > 0.95$ while carrying no recoverable $\widehat{\Mat{W}}$ and learning
Jacobians inconsistent with the wiring; the edge-ablation certificate, which they
structurally cannot pass and the GNN passes at $r \geq 0.99$ without retraining;
tolerance to $+400\%$ random and $+661\%$ spatially plausible false edges
($\Rw \approx 0.96$ at $5\times$ the true edge count); the FlyWire-scale runs at
$50{,}412$ neurons and $9.6$M candidate edges; and $\Mat{W}$ and $\tau$ recovered
on par with the oracle at $\sigma > 0$ ($\Rw = 0.99$ for both, full-sample
$\Rt = 0.99$ against $1.00$). Two results we do withdraw: $V^{\mathrm{rest}}$
recovery and the SIREN stimulus claim. We take the concessions below to be
bounded in that way.

\paragraph{1. Identifiability analysis internally inconsistent.}
\emph{Eq.~(11) is a count, not a kernel.} $308{,}160 = \sum_{i,\alpha}
(k_{i\alpha}-1)$ is correct as computed: it counts sum-zero directions within
same-type presynaptic groups, a structural property of the connectome. Calling
that set a kernel was the error. Same-type neurons deliver near-identical,
time-shifted copies of one signal, so raising one input and lowering a sibling
leaves the summed drive almost unchanged; but time-shifted signals are
correlated, not identical, so those directions carry small \emph{nonzero}
singular values.
\phantomsection\label{link1}
\emph{The two percentages are two points on one curve.} They answer different
questions at different tolerances on the same spectrum, which the submission
never said. Ranking weight directions by $\epsilon = \lambda_k/\lambda_{\max}$
and counting the share of edges degenerate at each tolerance, the numerical
$6.4\%$ is the value at $\epsilon = 10^{-22}$ and Eq.~(11)'s $71\%$ is reached at
$\epsilon = 2.5\times10^{-6}$, read off the same $\sigma = 0$ curve rather than
chosen. A strict tolerance counts only edges the data cannot constrain at all; a
looser one also counts weakly constrained columnar ones.
\phantomsection\label{link2}
\emph{Why process noise helps, now measured.} Appx.~C writes the noise-free ODE
as a per-neuron linear system $A_i \theta_i = b_i$, one row per timestep and one
column per unknown of neuron $i$. We take the singular values of $A_i$ by SVD and
rank each weight direction by $\epsilon = (\sigma_k/\sigma_1)^2$; a direction
below a fixed tolerance is one the activity does not constrain, and an edge
counts as degenerate once a direction supporting it falls below it. No weights
are fitted and no solver runs, so this measures the conditioning of the inverse
problem, not the behaviour of our GNN.

\begin{center}
{\footnotesize
\begin{tabular}{lcc}
\toprule
process noise $\sigma$ & degenerate $W$ null (\%) & degenerate $W$ sloppy (\%) \\
\midrule
$0$    & $6.35 \pm 0.02$ & $4.01 \pm 0.36$ \\
$0.05$ & $2.44 \pm 0.01$ & $0.32 \pm 0.00$ \\
$0.5$  & $0.00 \pm 0.00$ & $0.00 \pm 0.00$ \\
\bottomrule
\end{tabular}}

{\footnotesize Tolerances $\epsilon \leq 10^{-22}$ (null) and $\leq 10^{-12}$
(sloppy); mean\,$\pm$\,SD over five folds. The horizon is fixed across $\sigma$
and $\epsilon$ is scale-invariant, so the trend is not an artefact of record
length or of noise inflating variance.}
\end{center}

Two competing readings are excluded. That noise merely smooths the GNN's loss
landscape is ruled out by the oracle, which has no learned functions to
regularise yet improves identically ($\Rw: 0.96 \to 0.99 \to 1.00$). That noise
lifts the spectrum by a numerical floor is ruled out by column equilibration,
which fixes the sum of squared singular values, so noise can only redistribute a
fixed total by decorrelating columns. We claim only what the table supports:
process noise improves the conditioning. Whether it acts selectively on the
columnar directions or lifts the spectrum generally the degenerate-edge fraction
cannot say, and we will report that rather than assert it.

\textbf{What we'll change.} State both thresholds as $\epsilon \leq 10^{-22}$ and
$\leq 10^{-12}$, which reproduce the reported $\approx\!10\%$ where the printed
$64000\times$ factor does not; replace ``kernel'' with $\ker_\epsilon$ at that
tolerance, keeping $308{,}160$ as an exact count; split Appx.~C into C.1
(identifiability) and C.2 (dynamical equivalence at an explicit tolerance),
promoting C.2 as the analytical backbone of the prediction $\neq$ mechanism
claim; and replace the asserted noise interpretation, in the appendix and in
every claim citing it, with the measurement above.

\paragraph{2. The setting is strongly model-matched.}
Partially accepted. Of the nine listed items, three are properties of the setup
rather than the method, since a measured connectome, a known stimulus and a
blank interval are all available in a real experiment and the paper already
degrades the first two; two are inaccurate as stated, since no sign convention
is imposed and only the aggregation is additive; and four are genuine:
first-order dynamics, direct voltage observation, matched $\Delta t$, and the
monotonicity priors. We tested three of the four. The GNN holds within $0.023$
$\Rw$ of the oracle out to $\Delta t_{\mathrm{obs}} = 40$~ms and falls to
$0.315$ at $100$~ms; an unobserved adaptation current costs $0.078$ at
$g_a = 0.3$; the monotonicity priors are not load-bearing. Numbers in our reply
to Reviewer 1bPK, Q1.

\textbf{What we'll change.} Scope the title, abstract and contributions to
model-matched in-silico feasibility, and replace ``without assuming the form of
the dynamical equation'' with the class actually assumed: first-order dynamics
with additive pairwise aggregation over the supplied graph, agnostic only to the
scalar nonlinearities.

\paragraph{3. Biological and practical relevance.}
Conceded on the premise: no dataset pairs synapse-resolution connectivity with
dense activity in one animal. Measurement noise remains the bottleneck,
$\Rw = 0.63$ and $\Rt = 0.28$ at $\gamma = 0.1$, $0.38$ and $0.08$ at
$\gamma = 0.2$, with the oracle failing alongside, which places the limit in the
inverse problem rather than in our architecture. On the recurrent circuit, the
\emph{Drosophila} CX run gives $\Rw = 0.54$ noise-free against Flyvis's $0.89$,
rising to $1.00$ at $\sigma = 0.5$ (Vzfg, Q1). The calcium observation model,
now run, is the harder result: fitted on the indicator trace the GNN recovers no
weight structure at any degradation (Vzfg, Q2). We report it as the negative
result it is.

\textbf{What we'll change.} The measurement-noise result bounds the claim rather
than adding to it, so it moves out of Limitations into the contribution list and
the abstract, which will state recovery as demonstrated under process noise and
\emph{not} at realistic measurement noise. We will also present the optogenetics
analogy as an untested hypothesis rather than a prescription.

\paragraph{4. ``Mechanism recovery'' is ambiguous.}
Accepted; the paper uses one phrase for four claims of decreasing strength.
Exact parameter equality we do not claim. Recovery up to an equivalence class we
do claim, for $\sigma > 0$. Agreement of the local vector field we now measure:
$R^2 = 0.970$ against the true Jacobian, with the oracle at $0.977$. Counterfactual
competence under edge ablation is necessary but not sufficient, since it does not
establish uniqueness. The third is weaker than the second, and we say so: a
Jacobian is $\Mat{W}$ rescaled by $1/\tau_i$ and by each source's active
fraction, so agreement there cannot separate a recovered connectome from a
dynamically equivalent one.

\textbf{What we'll change.} State the four senses explicitly in Sect.~1 and attach
each claim to one of them; fix the excluded set a priori rather than by residual,
and report full-sample metrics for every table (Reviewer 1bPK, Q2).

\paragraph{5. Methodological novelty overstated.}
Partially accepted. The decomposition is inherited from \textbf{Allier et al.\
[2026]}; \textbf{AMAG [2023]} refines an adjacency it initialises rather than
inferring one; \textbf{Yoon [2025]} infers a
$100$-neuron ring from $19$ observed cells, recovering a smooth weight profile
rather than individually free weights. The degeneracy itself is older still,
Prinz, Bucher and Marder [2004] and Golowasch, Goldman, Abbott and Marder
[2002], with Das and Fiete drawing the consequence for inference. The defensible
contribution is the connectome-scale characterisation, the identifiability
analysis, the degradation battery and the edge ablation.

\textbf{What we'll change.} Reframe the contribution as a large-scale benchmark of
these ideas rather than their origin, and add a Sect.~2 paragraph separating
forward connectome models (Shiu) from fits to activity (BrainTrace, Mi,
Pospisil, Lappalainen).

\paragraph{6. Agentic hyperparameter search.}
Accepted in full: no matched-compute baseline against random or Bayesian search,
and it optimises $\Rw$ against ground truth, which is unavailable in the intended
application.

\textbf{What we'll change.} Demote it from contribution~(3) to a methods note,
replace the raw logs with a curated summary, remove the self-assessments, and
retain only the outcome: one configuration used unchanged everywhere.

\paragraph{Reporting issues.}
$N$ reconciled: $13{,}697$ neurons have at least one incoming edge and $44$ have
none, so Appx.~C counts the former. The $\pm 0.00$ entries are rounding, not
under-dispersion.

\textbf{What we'll change.} Report three decimals throughout, spell out SIREN at
first use, add unfiltered metrics to all four tables, and reset Figs.~1 and~2 at
main-text font size, splitting Fig.~1's extraction panel into the section that
describes it.

\bigskip
\hrule
\bigskip

\section*{Reviewer 1bPK READY TO SUBMIT ?}
We thank the reviewer for the precise review. We address the reviewer's comment below.

\paragraph{Q1. How much of "recovery" depends on (a) the GNN's hypothesis class containing the true generator and (b) $\Delta t_{inference} = \Delta t_{simulation}$ ? } Notably we already measure the effect of temporal mismatch in Supp. Table 4, row 1/5 frames.  And the limitation in acknowledged \emph{'Inference and simulation share $\Delta t = 20\,\text{ms}$, so the piecewise-linear update is exact rather than approximate'}. We propose to add new temporal mismatch results (see table below) and better discuss the recovery degression that follows.

\begin{center}
{\footnotesize
\begin{tabular}{cccc}
\toprule
$\Delta t_{simulation}$. & $\Delta t_{inference}$ & GNN $\Rw$ & Known ODE $\Rw$ \\
\midrule
$20$~ms & $20$~ms  & $0.970$    & $0.955$ \\
$10$~ms & $20$~ms  & $0.864$    & $0.873$ \\
$4$~ms  & $20$~ms  & $0.822$    & $0.844$ \\
$2$~ms  & $20$~ms  & $0.812$    & $0.835$ \\
\midrule
$2$~ms  & $40$~ms  & $0.720$    & $0.711$ \\
$2$~ms  & $100$~ms & $0.315$    & $0.449$ \\
\bottomrule
\end{tabular}}
\end{center}

\phantomsection\label{link3}

About model-class match, we consider that the experiments in section 4, with false-positive connectivity to be sound model mismatch tests.
Agreed with the reviewer, we will discuss two new model-class match experiments. 
First, we test the current GNN, e.g.\ by adding a dynamical term outside the additive-pairwise class. 
In this new experiment, we add an unobserved per-neuron adaptation current:
\[
  \tau_i \dot v_i = -v_i + V_i^{\mathrm{rest}}
  + \sum_j W_{ij}\,\mathrm{ReLU}(v_j) - g_a c_i + I_i,
  \qquad
  \dot c_i = (v_i - c_i)/\tau_a ,
  \qquad
  \tau_a = 200~\mathrm{ms} \gg \tau
\]
This violates model-class match twice. First, the observed system is no longer
first order. Second, $c_i$ is a state variable never seen during training; the
GNN is unchanged.
 
\begin{center}
{\footnotesize
\begin{tabular}{lccc}
\toprule
$g_a$ & GNN $\Rw$ & GNN $\Rv$ & Known-ODE $\Rw$ \\
\midrule
$0$   & $0.987$ & $0.97\,(0.63)\,[5.0]$   & $0.986$ \\
$0.1$ & $0.974$ & $0.80\,(0.62)\,[5.7]$   & $0.972$ \\
$0.3$ & $0.909$ & $0.54\,(-0.35)\,[27.7]$ & $0.923$ \\
\bottomrule
\end{tabular}}
 
\tabnote{$\Rv$: inlier (full-sample) [excl.\ \%]. $\Rt \geq 0.98$ in every row.}
\end{center}
 
Degradation is ordered and the oracle matches: $\widehat V^{\mathrm{rest}}$
absorbs the slow offset first, $\widehat W$ inherits its bias, $\tau$ (from the
slope of $f_\theta$) survives.
Here recovery does not degrade sharply. 
The reviewer also suggests testing the GNN against a conductance-based
simulation. The generator becomes
\begin{equation}
\tau_i\frac{dv_i(t)}{dt} = -v_i(t) + V_i^{\mathrm{rest}}
+ \sum_{j\in\mathcal{N}_i} W_{ij}\,\mathrm{ReLU}\!\big(v_j(t)\big)\,
  \big(E_{ij} - v_i(t)\big)
+ I_i(t) + \sigma\,\xi_i(t),
\label{eq:conductance_update}
\end{equation}
where $E_j$ is the reversal potential of the synapses made by neuron $j$, set by its
neurotransmitter. The synaptic current is therefore multiplicative in
the postsynaptic state, and the sign of the interaction is carried by
$E_j - v_i$ rather than by the sign of $W_{ij}$. This second experiment is running, we intuit that this mismatch can affect recovery.
To address conductance-based model, it will be necessary to modify the GNN form to:
\begin{equation}
\frac{d\hat{v}_i(t)}{dt}
= f_\theta\!\Big(v_i(t),\,\mathbf{a}_i,\,
\sum_{j\in\mathcal{N}_i} \hat{W}_{ij}^2\,
g_\phi\!\big(v_i(t), v_j(t),\mathbf{a}_i,\mathbf{a}_j\big),\,
I_i(t)\Big).
\label{eq:gnn_ct}
\end{equation}
We intend to address conductance-based simulation with this modified GNN form in future works though.


Additionally, we checked that the prior is not load-bearing. On Flyvis-217
blank50 at $\sigma = 0.05$, single seed: $\mu_1$ on gives $\Rw = 0.987$, $\mu_1$
off gives $0.988$. ($\mu_0$ is already $0$ in the published configuration, so
there is nothing to switch off there; a repeat of the $\mu_1$-on run also lands
at $0.988$, so the two are equal to within run-to-run variation.)
Note that there is no sign convention imposed, $\hat W_{ij}$ carries the sign and $g^2_{\phi}$ the magnitude.


\phantomsection\label{link4}
\paragraph{Q2. Please report the full-sample (no outlier filtering).}
We computed full-samples for all four tables (Tabs.~1, 2, Supp.~4, 7) and  we will report in the format: inlier (full-sample) [excluded \%].
Obviously, where outlier exclusion is small the two agree; where it is large the inlier value is not informative. 

\paragraph{Q3. Reconcile "on par with an oracle". }
Agreed with reviewer, we will restrict the claim to model noise $\sigma > 0$. Noise-free regime is where the oracle better performs.

\paragraph{Q4. Benchmark for agentic loop.}
We have no such baseline, no random and/or Bayesian search was run. 
So the claim that our agentic loop yields better configurations is
unsupported. We will keep the agentic loop in methods to explain how we settle one training configuration used everywhere (Supp.~6).

\phantomsection\label{link8}
\paragraph{Q5.&Q6. Previous work.}
\textbf{Allier et al.\ [2026]} preprint appeared February 2026, three months
before the deadline. It is a different system and a different formulation. 
The tested neural assemblies are random with no anatomy, whose weights carry no columnar degeneracy,
Their GNN formulation is additively separable:
\begin{equation}
\frac{d\hat{v}_i(t)}{dt} =\phi^*\left(\mathbf{a_i}, v_i\right) + \Omega_i^*\left(t\right) \sum_{j\in\mathcal{N}_i} \Mat{W}_{ij}\psi^*\left(\mathbf{a_i},\mathbf{a_j},v_j\right).
\label{eqn:GNN}
\end{equation}
where our GNN formulation nests them to 
\begin{equation}
\frac{d\hat{v}_i(t)}{dt}
= f_\theta\!\Big(v_i(t),\,\mathbf{a}_i,\,
\sum_{j\in\mathcal{N}_i} \hat{W}_{ij}\,
g_\phi\!\big(v_j(t),\mathbf{a}_j\big)^2,\,
I_i(t)\Big).
\label{eq:gnn_ct}
\end{equation}
Our GNN formulation is a more general form, that does not map directly to the simulator ODE as their. Conversely their form features a modulation term $\Omega_i^*\left(t\right)$ that our lacks. 
To relate this to Q1, we will discuss different form of the GNN as being complementary to address different class of systems. Ultimately, we envision a general form combining all discussed that remains to be tested:
\begin{equation}
\frac{d\hat{v}_i(t)}{dt}
= f_\theta\!\Big(v_i(t),\,\mathbf{a}_i,\,
\Omega_i\left(t\right) \sum_{j\in\mathcal{N}_i} \hat{W}_{ij}^2\,
g_\phi\!\big(v_i(t), v_j(t),\mathbf{a}_i, \mathbf{a}_j\big),\,
I_i(t)\Big).
\label{eq:gnn_ct}
\end{equation}
\phantomsection\label{link9}
\\
\textbf{Shiu et al.\ [2024]} is a forward model rather than an inversion: per-edge
weights are taken from the connectome, their sign from the predicted
neurotransmitter, and a single global scalar is calibrated once against one
physiological reference point. On that basis it matches $91\%$ of $164$
empirically tested predictions.
That bears on our motivating sentence, ``top-down, hand-specified large-scale
simulations have struggled to be predictive of circuit function''. Shiu et al.\
is bottom-up and connectome-derived, so it does not contradict the sentence as
written, but the sentence reads as a claim about scale and we will qualify it:
what has proven hard is top-down specification, not scale, and a connectome
supplies exactly the constraint that was missing. That strengthens rather than
weakens our premise. The inverse direction, recovering the parameters rather
than assuming them, is what remains open.
\phantomsection\label{link10}
\\
\textbf{BrainTrace [2026]} is a learning system for spiking networks;
its whole-brain \emph{Drosophila} fit poses the same inverse problem as our
Known ODE oracle: form fixed (the Shiu LIF), connectome supplying the support of
$\Mat{W}$, weights optimised to match recordings.  Their setting is
harder though. The loss is on \emph{neuropil-level} rates, so each residual
constrains a sum over thousands of neurons, and supervision at $1.2$~Hz is
$42\times$ our $\Delta t = 20$~ms. Our sweep reaches only $5\Delta t$,
where the GNN already falls. We will add this comparison to the paper.

\paragraph{W4. The measurement-noise result bounds the contribution}
Agreed. This will not stay in limitations. We will temper the significance framing.

\paragraph{W7. Effect of model noise.}
Appx.~C writes the noise-free ODE
as a per-neuron linear system $A_i \theta_i = b_i$, one row per timestep and one
column per unknown of neuron $i$. We will add a new analysis, take the singular values of $A_i$ by SVD and
rank each weight direction by $\epsilon = (\sigma_k/\sigma_1)^2$; a direction
below a fixed tolerance is one the activity does not constrain, and an edge
counts as degenerate once a direction supporting it falls below it. No weights
are fitted and no solver runs, so this measures the conditioning of the inverse
problem, not the behaviour of our GNN.
\begin{center}
{\footnotesize
\begin{tabular}{lcc}
\toprule
process noise $\sigma$ & degenerate $W$ null (\%) & degenerate $W$ sloppy (\%) \\
\midrule
$0$    & $6.35 \pm 0.02$ & $4.01 \pm 0.36$ \\
$0.05$ & $2.44 \pm 0.01$ & $0.32 \pm 0.00$ \\
$0.5$  & $0.00 \pm 0.00$ & $0.00 \pm 0.00$ \\
\bottomrule
\end{tabular}}

{\footnotesize Tolerances $\epsilon \leq 10^{-22}$ (null) and $\leq 10^{-12}$
(sloppy); mean\,$\pm$\,SD over five folds. The horizon is fixed across $\sigma$
and $\epsilon$ is scale-invariant, so the trend is not an artefact of record
length or of noise inflating variance.}
\end{center}
The possibility that noise lifts the spectrum by a numerical floor is ruled out by column equilibration,
which fixes the sum of squared singular values, so noise can only redistribute a
fixed total by decorrelating columns. We claim only what the table supports:
process noise improves the conditioning. Whether it acts selectively on the
columnar directions we cannot say. We will report these new measurements.

Optogenetics is an
untested analogy, so we will claim only that noise and perturbation
help in simulation. 
\paragraph{W8. SIREN} the SIREN rollout is scored on training frames and the
stimulus is largely present in the photoreceptors, so ``recovers unknown visual
stimuli'' comes out of the abstract, the experiment staying in the appendix as
partial recovery.

\paragraph{W9. Statistical-variability reporting}
We checked, \emph{$\pm 0.00$.} is rounding. At four decimals the GNN
folds at $\sigma = 0.05$ span $0.9829$ to $0.9886$ (SD $0.0021$) against $0.0002$
We will report three decimals. We will fix N = 13,741 in Appx. C.


\bigskip
\hrule
\bigskip

 
\section*{Reviewer Vzfg [DEPRECATED] }
We thank the reviewer for the precise review. We address the reviewer’s comment below.

\phantomsection\label{link8}


\paragraph{Q1. Recurrent benchmarks.}
Flyvis is indeed a feedforward and the paper measures it: a
stimulus-only MLP never seeing the connectome predicts every voltage at
$r = 0.97$ from $16$ past frames (Appx.~H.2), and Limitation~5 states that
stronger recurrent regimes are untested. The reviewer therefore ask to test more challenging recurrent settings. Note, \textbf{Allier et al.\ [2026]} already test GNN on recurrent setting but the tested neural assemblies are random with no anatomy. Here, we propose to add results on a new experiment: test the GNN on the \emph{Drosophila} central-complex head-direction circuit: $156$ neurons, $7$ types, $10{,}263$ hemibrain edges, $42\%$ density, $180\times$
Flyvis's.  In our simulation, angular velocity, an Ornstein--Uhlenbeck process, drives $20$ PEN\_a neurons on every frame ($2$ channels $\times$ $10$ neurons, weights in $\{0,1\}$). Visual landmark cues reach the $46$ EPG neurons over the first $100$ out of $10{,}000$ frames. The remaining $90$ neurons receive no input at any time. The circuit was not optimized for the heading task. As such, this is not a functional central-complex head-direction circuit but a recurrent and anatomically sound system that can complement Flyvis's experiment. As in the paper the binary adjacency matrix is given and the signs are not. Each neuron's E/I identity
has to be recovered, under only a weak Dale-consistency penalty that never sees
a ground-truth sign. The deterministic model noise-free regime ($\sigma=0$)is again the harder case, and model noise is again what breaks the degenerate directions and allows the GNN to match the Known-ODE upper bound, on every parameter we can score here (no bias voltage, no different cell types). As suggest by the reviewer we propose to add this new experiment to demonstrate application to a recurrent settings.                              


\begin{center}
{\scriptsize
\setlength{\tabcolsep}{5pt}
\begin{tabular}{ll cc ccc}
\toprule
model & $\sigma$ & one-step $r$ & rollout $r$ & $\Rw$ & $\hat\tau$ error & E/I sign \\
\midrule
GNN       & $0$    & $0.999$ & $1.000$ & $-0.43$ & $12.0\%$ & $84.1\%$ \\
          & $0.05$ & $1.000$ & $1.000$ & $0.65$  & $0.3\%$  & $99.6\%$ \\
          & $0.5$  & $0.999$ & $0.999$ & $1.00$  & $0.2\%$  & $100\%$ \\
\midrule
Known-ODE & $0$    & $0.996$ & $0.997$ & $-23.9$ & $45.6\%$ & $64.4\%$ \\
          & $0.05$ & $1.000$ & $1.000$ & $1.00$  & $0.0\%$  & $100\%$ \\
          & $0.5$  & $1.000$ & $1.000$ & $1.00$  & $0.0\%$  & $100\%$ \\
\bottomrule
\end{tabular}}

\end{center}


\paragraph{Q2. Calcium-like observation model.}
Accepted, and it is the change that most improves transfer to real recordings.
The submission's degradations never compose: the $1/5$-frames row
($\Rw = 0.30$) is subsampling alone and the $\gamma = 0.1$ row ($0.63$) is
additive noise on the voltage, and neither is an indicator model.

We built one with four parts, on the same trajectories at process noise
$\sigma = 0.05$ and the paper's $\Delta t = 20$~ms: a GCaMP6f kernel, saturation
past $K_d$ set at $1.57\times$ the trace SD, a $500$~ms exposure ($2$~Hz, a box
average over $25$ steps, which is the anti-aliasing filter a slower camera
applies), and shot noise of SD $\gamma$ relative to the trace. Fitting the
fluorescence directly recovers nothing under any of them, the kernel alone
included, so we report instead what an experimenter actually fits: the trace
deconvolved back to voltage by a Tikhonov inverse in Fourier space, $\lambda$
retuned at each noise level by maximising $r(\dot v)$ so every row is the best
inversion available and not a fixed filter. The grid below isolates the part that
turns out to decide the outcome, photon noise on the kernel observable; both arms
train on the paper's schedule, $13{,}741$ neurons.

How well the inversion recovers the voltage, before any fitting:

\begin{center}
{\scriptsize
\setlength{\tabcolsep}{5pt}
\begin{tabular}{lccc cc cc}
\toprule
& & & & \multicolumn{2}{c}{train} & \multicolumn{2}{c}{test} \\
\cmidrule(lr){5-6}\cmidrule(lr){7-8}
$\gamma/\mathrm{SD}(C)$ & $\gamma$ & SNR (dB) & $\lambda$ & $r(v)$ & $r(\dot v)$ & $r(v)$ & $r(\dot v)$ \\
\midrule
$0$    & $0$       & $\infty$ & $10^{-6}$ & \good{$0.9998$} & \good{$0.9933$} & \good{$0.9993$} & \good{$0.9747$} \\
$0.01$ & $0.0064$  & $40.0$   & $10^{-3}$ & \good{$0.9797$} & $0.4353$        & \good{$0.9829$} & $0.5246$ \\
$0.03$ & $0.0193$  & $30.5$   & $10^{-2}$ & \good{$0.9695$} & \bad{$0.2803$}  & \good{$0.9726$} & $0.3470$ \\
$0.10$ & $0.0644$  & $20.0$   & $10^{-1}$ & \good{$0.9562$} & \bad{$0.1719$}  & \good{$0.9582$} & \bad{$0.2144$} \\
\bottomrule
\end{tabular}}

{\footnotesize Per-neuron Pearson $r$ between the deconvolved estimate and the
true voltage, and between their time derivatives. $\lambda$ is retuned at each
level by maximising $r(\dot v)$, so every row is the best inversion available
there rather than a fixed filter. The trace is recovered essentially perfectly
throughout; its derivative, which is what $\Mat{W}$ is read from, is not.}
\end{center}

The fits that follow, on those four deconvolved recordings:
\phantomsection\label{link5}
\begin{center}
{\scriptsize
\setlength{\tabcolsep}{3.5pt}
\begin{tabular}{ll cc c cc c}
\toprule
model & $\gamma/\mathrm{SD}$ (SNR) & one-step $r$ & rollout $r$ & $\Rw$ & $\Rt$ & $\Rv$ & cluster \\
\midrule
GNN       & $0$ ($\infty$)  & \good{$0.974$} & \good{$0.992$} & \good{$0.969$} & \good{$0.96$} $(0.95)$ $[0.2]$ & $0.85$ $(0.54)$ $[5.5]$ & $0.87$ \\
          & $0.01$ ($40$)   & $0.656$ & $0.667$ & \bad{$0.183$}  & $0.37$ $(0.08)$ $[9.8]$          & $0.46$ $(-1.4)$ $[37.9]$  & $0.76$ \\
          & $0.03$ ($30$)   & $0.550$ & $0.403$ & \bad{$-1.270$} & $0.61$ $(-5.7)$ $[81.9]$         & $0.66$ $(-19.6)$ $[71.3]$ & $0.73$ \\
          & $0.10$ ($20$)   & $0.414$ & \bad{$0.115$} & \bad{$-0.794$} & \textcolor{gray}{---} $(-35.3)$ $[100]$ & $0.54$ $(-16.2)$ $[67.3]$ & $0.68$ \\
\midrule
Known-ODE$^{\dagger}$ & $0$ ($\infty$)  & \good{$0.972$} & \good{$0.991$} & \good{$0.915$} & \good{$0.95$} $(0.95)$ $[0.1]$ & $0.80$ $(0.31)$ $[10.2]$ & $0.90$ \\
          & $0.01$ ($40$)   & $0.643$ & $0.654$ & $0.383$        & \bad{$0.30$} $(0.25)$ $[2.2]$   & $0.64$ $(-2.1)$ $[40.8]$  & $0.86$ \\
          & $0.03$ ($30$)   & $0.536$ & $0.402$ & \bad{$0.181$}  & \bad{$-0.08$} $(-5.5)$ $[71.4]$ & $0.63$ $(-4.6)$ $[49.4]$  & $0.82$ \\
          & $0.10$ ($20$)   & $0.400$ & \bad{$0.015$} & \bad{$-0.060$} & \bad{$-8.67$} $(-38.3)$ $[95.0]$ & $0.56$ $(-7.9)$ $[57.6]$ & $0.74$ \\
\bottomrule
\end{tabular}}

{\footnotesize Tab.~1's columns and fields. $\Rt$ and $\Rv$ are reported inlier
(full-sample) [excluded \%] over $13{,}741$ neurons; where the excluded fraction
is large the inlier value is not informative, and \textcolor{gray}{---} marks a
cell with no inliers left. Colouring follows the paper, \good{green} above $0.9$
and \bad{orange} below $0.3$. $^{\dagger}$The oracle's parameter columns are read
at $32$k iterations rather than at the end: on this observable it converges there
and then drifts monotonically away, $\Rw = 0.915 \to 0.481$ at $\gamma = 0$ over
the remaining $1.5$M iterations, while its activity prediction stays intact, so
the endpoint would report a diverged fit and not a worse one. Everything else,
including both prediction and clustering columns, is end-of-training. The GNN
does not drift at $\gamma = 0$ but does at the two highest noise levels, where
its best $\Rw$ over training is $-0.011$ and $-0.010$. The $\gamma = 0$ row is a
control: with a known kernel and no noise $C = K * v$ is exactly invertible, so
it checks that the observation model and its inverse agree and says nothing about
calcium imaging.}
\end{center}

Deconvolution works, and that is worth stating plainly: given a clean trace it
returns the weights in full, $\Rw = 0.97$ against the $0.984$ voltage control,
with $\tau$, $V^{\mathrm{rest}}$ and cell-type clustering at the levels the paper
reports on voltage. The obstacle is not the indicator kernel, which is
invertible, but the photon noise that inverting it amplifies. Two-photon GCaMP
sits at roughly $20$ to $40$~dB, and across that whole range recovery is gone:
even reading each arm where it does best, the most either reaches at
$40$~dB is $0.38$, and by $20$~dB neither leaves its initialisation. The
columns show why. The trace stays faithful everywhere ($r \geq 0.96$) while
its derivative collapses, and $\Rw$ follows the derivative rather than the trace:
process noise is what makes $\Mat{W}$ identifiable and it is white, so the
identifying signal sits at high frequency, exactly the band the indicator
attenuates and the band inversion cannot restore without amplifying whatever
noise has filled it. A high $r$ on a calcium trace is therefore consistent with
the weights being entirely unrecoverable.

We also correct an error of our own. An earlier version of this reply carried
deconvolved fits at $\Rw = 0.07$ and read them as evidence against
deconvolution. Those runs used a derivative-penalising Tikhonov prior carried
over from a noisy setting; on a clean trace it suppressed about $70\%$ of the
derivative amplitude while leaving $r(v) = 0.98$, destroying the very quantity
being measured. That number characterised our preprocessing, not the observation
model, and we withdraw it.

\paragraph{Q3. Out-of-distribution test.}
Our evidence for ``black boxes predict but do not recover'' is indirect; this
test makes it direct and needs no retraining. Only the drive $I_i(t)$ changes:
from a shared initial state each frozen checkpoint is rolled out on its own
output for $1000$ steps ($20$~s), scored over the whole trajectory against
ground truth integrated from Eq.~1 under that drive. A model that recovered the
mechanism should hold in states it never saw.
\phantomsection\label{link6}
\begin{center}
{\footnotesize
\begin{tabular}{lcc}
\toprule
stimulus & GNN roll.\ $r$ & Known-ODE roll.\ $r$ \\
\midrule
naturalistic held-out (ID reference) & $1.0000$ & $1.0000$ \\
full-field white noise               & $0.9986$ & $0.9992$ \\
current injection, $5\%$ of neurons, $5\sigma_I$  & $0.9895$ & $0.9998$ \\
current injection, $5\%$ of neurons, $20\sigma_I$ & $0.9748$ & $0.9985$ \\
\bottomrule
\end{tabular}}
\end{center}

The two stimuli move progressively off the manifold. White noise is one time
series broadcast to every photoreceptor, and it barely moves either model. That
is informative: an unnatural stimulus is not an unfamiliar state, since the drive
still enters through the photoreceptors and the circuit visits states it was
fitted on. Injection adds a constant to a random $5\%$ of \emph{all} neurons,
most beyond the reach of any stimulus, forcing voltages the visual pathway cannot
produce. There a gap opens and grows with amplitude, $0.0008$ at $1\sigma_I$,
$0.0103$ at $5\sigma_I$, $0.0237$ at $20\sigma_I$ ($\sigma_I$ = SD of the visual
drive), so $1\sigma_I$ is not yet out of distribution and sets the discriminating
scale. The oracle holds above $0.998$, knowing Eq.~1 and fitting only constants,
where the GNN's $f_\theta, g_\phi$ are evaluated outside their training range.

\emph{What this settles.} It bounds our own off-manifold fidelity but not the
comparison asked for: the gap is between our model and an oracle knowing the
equation, not between a mechanistic model and one that only predicts. The
recurrent-MLP and EED checkpoints are with an absent co-author; that column needs
rollouts, not retraining, and will be filled during the discussion.

\paragraph{Q4. GNN and Known-ODE Jacobians.}
Agreed: Supp.~Fig.~18 uses a Jacobian mismatch against the MLP and never
applies the same test to our own model.

At each of $200$ held-out frames we compute
$J_{ij} = \partial \dot v_i / \partial v_j$ for every model. Differentiating
Eq.~1 gives $J_{ij} = W_{ij}\,\mathrm{ReLU}'(v_j)/\tau_i$, $J_{ii} = -1/\tau_i$;
the Known-ODE has the same form in $\widehat W, \hat\tau$, and the GNN is
differentiated by autograd through $f_\theta, g_\phi$, since the ReLU kinks make
finite differences unusable. All models see the same states, on the noise-free
blank50 run of Supp.~Fig.~18.
\phantomsection\label{link7}
\begin{center}
{\footnotesize
\begin{tabular}{llcccc}
\toprule
block & model & $R^2$ vs GT $J$ & Pearson $r$ & $|J|$ off graph & sign agr. \\
\midrule
off-diagonal $J_{ij}$ & Known-ODE & $0.977$ & $0.988$ & $0$ by construction & $0.882$ \\
                      & GNN       & $0.970$ & $0.988$ & $0$ by construction & $0.878$ \\
\midrule
diagonal $J_{ii}$     & Known-ODE & $0.944$ & $0.975$ & n/a                 & $1.000$ \\
                      & GNN       & $0.865$ & $0.963$ & n/a                 & $1.000$ \\
\bottomrule
\end{tabular}}

{\footnotesize MLP row pending; checkpoints inaccessible (Q3).}
\end{center}

The GNN matches the true off-diagonal Jacobian at $R^2 = 0.970$, within $0.007$
of the oracle, so it has the right local dynamics, not merely plausible
parameters. Not an artefact of silent edges: $92.9\%$ carry a nonzero
ground-truth Jacobian and restricting to those moves $R^2$ by $<10^{-3}$. The
diagonal is weaker, $0.865$ against $0.944$: the oracle holds $-1/\tau_i$ as a
parameter, the GNN expresses the leak through $\partial f_\theta / \partial v$.

A good Jacobian is not a recovered connectome. $J$ is $\Mat{W}$ rescaled twice:
by $1/\tau_i$, and, over frames, by the fraction in which neuron $j$ sits above
zero. Those rescalings change which edges dominate, which is why $J$ scores
higher than $\Mat{W}$: the GNN reaches $0.970$ against $\Rw = 0.913$, the oracle
$0.977$ against $0.964$, where the whole gap is the rescaling ($0.964 \to 0.969$
after $1/\tau_i$, $\to 0.977$ after the activity factor) with no change to the
fit. Agreement with the vector field therefore cannot separate a recovered
connectome from a dynamically equivalent one, our argument about BrainTrace to
Reviewer 1bPK.

Two entries carry less than they appear to: the off-graph fraction is zero by
construction for both our models, reporting the hypothesis class rather than what
was learned, and sign agreement reduces over frames to sign recovery of
$\Mat{W}$, the oracle at the same $0.882$.

\textbf{What we'll change.} Add these panels and metrics beside Supp.~Fig.~18,
with the MLP row once its checkpoints return, and state what this measurement
forces: recovering the vector field is weaker than recovering the connectome,
beside the Appx.~C degeneracy argument we so far apply only to other people's
fits.

\paragraph{Minor comments.}
Both agreed. The Appx.~F.2 mapping from $f_\theta, g_\phi$ to
$\hat\tau, \widehat V^{\mathrm{rest}}, \widehat W$ is what makes the GNN
interpretable rather than another black box, and we will move it to the main
text. We will also reset Figs.~1 and~2 at main-text size, splitting Fig.~1's
extraction panel out into the section that will describe it.

\bigskip
\hrule
\bigskip

\section*{Reviewer PVbi [DEPRECATED] }
We thank the reviewer. Two of the five weaknesses concern framing rather than
results and we accept both: the decomposition is inherited rather than new, and
``mechanism recovery'' covers several claims that are not equivalent. We answer
in the order of the letter, and where a point recurs across
reviews we answer in the same words: three inconsistent replies would only
confirm the imprecision you identify.

\paragraph{W1. Novelty against AMAG [2023] and Yoon [2025].}
Accepted that the decomposition is not new. The two cited works differ from
ours, and from each other, in what they infer and at what scale.
\phantomsection\label{link11}
\textbf{AMAG [2023].} It forecasts field potentials and $\mu$ECoG from macaque
electrode arrays, learning its adjacency adaptively from a correlation-matrix
initialisation ``to stabilize the learning process''. On synthetic linear systems
it does recover a known adjacency ($F_1 = 0.88$ to $0.93$, one-step forecasting),
but the units are electrode channels rather than identified neurons, and the real
recordings carry no ground-truth connectivity: validation is by channel spatial
proximity. Its own finding that recovery tracks the training
objective, $F_1$ falling to $0.63$ under reconstruction, is the dissociation we
report between fit quality and structure recovery.
\phantomsection\label{link12}
\textbf{Yoon [2025].} This work does infer connectivity: on a $100$-neuron ring
attractor with auxiliary nodes for unobserved cells, and on real head-direction
data where $19$ recorded cells sit in a ring of $N = 100$. Their
$\widehat{\Mat{W}}$ is fully free, predicted per edge by an MLP over learned
embeddings and scored on the whole matrix as
$\|\Mat{W}-\widehat{\Mat{W}}\|_F/\|\Mat{W}\|_F$. The target is what differs: their
ground-truth ring is rotationally symmetric, each row a phase-shifted copy of one
profile, so it holds $O(N)$ degrees of freedom rather than $N^2$, and the
evaluation uses that, aligning rows to a common phase and fitting one global
scale first. That is the same gain degeneracy our $\mu_2$ anchor addresses (W5).
Flyvis-217 is convolutional too, $434{,}112$ edges tiled from $2{,}355$
templates, and that symmetry is what Appx.~C measures; the differences are that
we do not exploit it in evaluation, and that our FlyWire hybrids break it with
column-resolved wiring, $\Rw = 0.94$ on $327{,}358$ measured edges and $0.97$ on
$1{,}266{,}378$ across $50{,}412$ neurons. The failure
they report in other methods, ``connections to neurons that are either
unconnected or only weakly connected'' driven by ``strong correlations in neural
activity'', is Das and Fiete's, separated from ours under W2.

\textbf{Ours.} A \emph{measured} connectome as a binary prior rather than a graph
to infer, at $13{,}741$ to $50{,}412$ neurons, with per-neuron $\tau$ and
$V^{\mathrm{rest}}$ alongside the weights, and the degeneracy quantified rather
than described: Appx.~C measures which weight directions the activity leaves
unconstrained, and our meta-review reply measures how noise removes them.

\emph{On the given graph.} The prior is not simply accepted; the paper
measures both directions. Inflated, the GNN zeroes the $80\%$ of false edges and
still recovers $\widehat{\Mat{W}}$ at $R^2 = 0.96$ at $5\times$ the true edge
count, which is edge inference within a supplied superset. Deleted, recovery
degrades at $-20\%$ ($\Rw = 0.78$) and breaks at $-50\%$ ($\Rw = 0.19$, rollout
$r = 0.56$). The method prunes but cannot add: it finds which supplied candidates
carry weight and fails when a true edge is absent. That is narrower than
inferring a graph, and we should have stated it as the boundary rather than
leaving the two ablations to be read separately.

\textbf{What we'll change.} State the delta explicitly in Sect.~2: a measured
connectome used as a prior rather than a graph to be inferred, at connectome
scale, with the degeneracy counted rather than described.

% W2 citations verified against PubMed 2026-07-26 (abstracts in papers/):
%  - Prinz, Bucher & Marder, Nat Neurosci 7(12):1345-52 (2004), doi:10.1038/nn1352.
%    "virtually indistinguishable network activity can arise from widely disparate
%    sets of underlying mechanisms" -- 20M+ simulated pyloric networks. Author
%    order confirmed.
%  - Golowasch, Goldman, Abbott & Marder, J Neurophysiol 87(2):1129-31 (2002),
%    doi:10.1152/jn.00412.2001. A model built from the MEANS of the one-spike
%    burster population "is not itself a one-spike burster, but rather fires three
%    action potentials per burst", because those sets "lie in a highly concave
%    region of parameter space that does not contain its mean". Author order
%    confirmed. No PDF: the publisher blocks automated download.
\phantomsection\label{link13}
\paragraph{W2. Das and Fiete [2020] precedent.}
Accepted, and the precedent runs further back than [3]. That many different
circuits produce indistinguishable activity is Prinz, Bucher and Marder's result
[Nat.\ Neurosci.\ 2004], and Golowasch, Goldman, Abbott and Marder
[J.\ Neurophysiol.\ 2002] showed those parameter sets occupy a non-convex region,
so their mean is not itself a valid model. That bears on our wording: claiming
recovery up to an equivalence class, reporting a representative member is not a
hedge but the only defensible form, since interpolating between members can leave
the class. Das and Fiete drew the consequence for inference, and we give them the
same answer as Reviewer Vzfg: their failure is a false edge where no synapse
exists, ours a true edge unidentifiable against its same-type neighbours, and we
will cite our process-noise result as a measured instance of their prescription
rather than an independent finding. We accept the reframing: what we add is not
the insight but its measurement at connectome scale, where the degenerate
directions can be counted rather than inferred from a failed fit.

\textbf{What we'll change.} Cite Prinz et al.\ and Das and Fiete at the point where
the identifiability claim is made, not only in related work, and describe the
contribution as a benchmark of these ideas rather than their origin.

\paragraph{W3. Flyvis is a simulator, not paired real data.}
Conceded, and your form of the point is the sharper one: Flyvis is a
hand-designed ODE with simplified neuron and synapse dynamics, so recovering it
demonstrates recovery of a simulator. No dataset pairs synapse-resolution
connectivity with dense activity in one animal, the premise of the work rather
than an oversight.

Two things bound the cost. The class is not arbitrary: the fly optic lobe is
non-spiking and graded-potential, so a first-order voltage model is the
appropriate form rather than a convenient one. And the hand-designed form is what
the new runs attack. We add an unobserved adaptation
current, neither first order in the observables nor decomposable into pairwise
messages, and we integrate the generator below the model's own step. Recovery
degrades in an ordered, measurable way rather than collapsing,
$V^{\mathrm{rest}}$ first, $W$ next, $\tau$ last (Reviewer 1bPK, Q1). How far the
result survives a wrong equation is the most this design supports without paired
data.

\textbf{What we'll change.} Retitle to ``\dots from simulated neural activity'',
scope the abstract and contributions to model-matched in-silico feasibility, and
name measurement noise rather than scale as the obstacle to real recordings.

\paragraph{W4. Senses of ``mechanism recovery''.}
Accepted; the paper uses one phrase for four claims of decreasing strength.
Exact parameter equality we do not claim. Recovery up to an equivalence class we
do claim, for $\sigma > 0$. Agreement of the local vector field we now measure:
$R^2 = 0.970$ against the true Jacobian, with the oracle at $0.977$.
Counterfactual competence under edge ablation is necessary but not sufficient,
since it does not establish uniqueness. The third is weaker than the second, and
we say so: a Jacobian is $\Mat{W}$ rescaled by $1/\tau_i$ and by each source's
active fraction, so agreement there cannot separate a recovered connectome from a
dynamically equivalent one. The second is also weaker than it sounds, for the
reason given under W2: the equivalent set is non-convex, so a representative is
all that can be reported.

\emph{On the filtering.} You are right that $\tau$ and $V^{\mathrm{rest}}$ are
extracted by local linearisation and reported after residual filtering, with
exclusions reaching $40.8\%$ in Tab.~2, $49.1\%$ under measurement noise and
$56.3\%$ under temporal subsampling. We have recomputed all four tables without
the filter. $\tau$ survives, full-sample $0.99$ against the oracle's $1.00$ at
$\sigma = 0.05$. $V^{\mathrm{rest}}$ does not: $0.76$ against $0.95$, and negative
in every degraded row. A high $R^2$ on the surviving subset does not license an
unqualified claim, and we withdraw it for $V^{\mathrm{rest}}$.

\emph{On the ablation.} Also accepted. The same random mask on two models sharing
a graph tests whether each uses that graph, not whether either has the right
weights on it: a necessary condition, presented as one.

\textbf{What we'll change.} State the four senses in Sect.~1 and attach every
claim to one; report full-sample metrics for all four tables, labelling a
parameter recovered only at full-sample $R^2 > 0.9$; and fix the excluded set a
priori, from zero in-degree ($44$ neurons) and neurons never crossing threshold,
so filtering is not defined by the residual it excuses.

\paragraph{W5. Clarity: relation to Flyvis, GNN against the oracle.}
Both points accepted. Flyvis obtained $\tau$, $V^{\mathrm{rest}}$ and $\Mat{W}$
by task optimisation on optical flow, never fitting activity; we recover the same
quantities from activity alone. The novelty is the inverse problem, not the
parameter types, and the paper should say so where it introduces them.

On the comparison: the oracle knows Eq.~1 and fits three parameter families,
while the GNN learns $f_\theta$, $g_\phi$, the embeddings $\Mat{a}_i$ and
$\widehat{\Mat{W}}$. You ask whether those learned functions are identifiable,
and the answer is only up to a gain: rescaling $g_\phi$ can be absorbed into
$\widehat{\Mat{W}}$ without changing the dynamics. That is why $\mu_2$ anchors
$g_\phi$ at a reference voltage, and why the Appx.~F.2 extraction exists at all.
The learned functions are not the parameters; the map between them is what makes
the GNN interpretable, and it is currently in an appendix.

\textbf{What we'll change.} State the Flyvis contrast where the parameters are
introduced, give the gain degeneracy and its fix beside the loss, and move the
extraction into the main text.

\paragraph{Q1.} SIREN is a sinusoidal representation network; we will spell it
out at first use.

\paragraph{Q2.} The $0.9$ green shading is a readability aid with no statistical
meaning; we will say so in the captions.

\section{Dynamical equivalence and non identifiability}

We recall the FlyVis ODE here for convenience.
\begin{equation}
    \tau_i\frac{dv_i(t)}{dt} = -v_i(t) + V^{\mathrm{rest}}_i
+ \sum_{j\in\mathcal{N}_i} W_{ij}\,\text{ReLU}\!\big(v_j(t)\big)
+ I_i(t) \label{eq:Flyvis-app-noise-free-ode}
\end{equation}
and collectively refer to the parameters $(\tau_i, V^{\mathrm{rest}}_i, W_{ij})$ as $\theta$.

\subsection{Definitions}

\paragraph{Degenerate/Null space/Non-identifiability}
Given activity traces $\tilde{v}_i(t)$, $t=1, \cdots, T$, the inverse problem of determining the parameters $\theta$ could admit multiple solutions in a strict, exact sense. For example, given a solution of the ODE $\dot{x} = - (a+b)x$, only the sum of the parameters $a, b$ can be inferred. Here we focus on continuous degeneracies.

\paragraph{Sloppy directions/approximate degeneracy}
These refer to directions in the parameter space that are poorly constrained by the observed activity traces, but not strictly degenerate. The Hessian matrix of the loss function evaluated at the true minimum is ill conditioned where the ratio of smallest to largest eigenvalues is $\ll 1$. We pick an arbitrary threshold of $1e-12$.

\paragraph{Dynamical equivalence}
By dynamical equivalence we mean directions in parameter space that leave the solution curves of an ODE unchanged. Either strictly unchanged like an exact parameter symmetry, or in an approximate manner. Concretely, for a fixed initial condition $v_i(0)$, if the continuous family of ODEs parameterized by $\theta + \lambda \eta$, where $\lambda \in \mathbb{R}$, for some fixed direction $\eta$ in parameter space produces exactly or approximately invariant $v_i(t), t > 0$ then the dynamics along the direction $\eta$ are equivalent (either exactly or approximately).

\subsection{}

We show in Appendix C we leverage the linearity of the FlyVis ODE in the parameters $\mathbf{\theta}_i$, to characterize the degenerate space as the kernel of a design matrix $\mathbf{A}_i$ in Suppl. Eq. 6.

%%% CONTINUE THIS AND SHOW THAT DYNAMICAL EQUIVALENCE = DEGENERACY

\end{document}
